% Generated draft: requires mathematical and source review.
\documentclass[11pt]{article}
\usepackage[margin=25mm]{geometry}
\usepackage{fontspec}
\setmainfont{Latin Modern Roman}
\usepackage{amsmath,amssymb,mathtools,unicode-math}
\setmathfont{Latin Modern Math}
\usepackage{enumitem,needspace,xurl,hyperref}
\setlength{\parindent}{0pt}
\setlength{\parskip}{5pt}
\setlength{\emergencystretch}{3em}
\newcommand{\N}{\mathbb{N}}
\newcommand{\Z}{\mathbb{Z}}
\newcommand{\Q}{\mathbb{Q}}
\newcommand{\R}{\mathbb{R}}
\newcommand{\C}{\mathbb{C}}
\newcommand{\F}{\mathbb{F}}
\newcommand{\A}{\mathbb{A}}
\newcommand{\PP}{\mathbb{P}}
\newcommand{\E}{\mathbb{E}}
\newcommand{\eps}{\varepsilon}
\newcommand{\dd}{\,\mathrm{d}}
\begin{document}
\section*{UnsolvedMath ID: 30002654}
\textbf{Source problem code:} catalogue entry
\subsection{Definitions and mathematical statement}
The supplied record describes a conjectural moduli and filling problem, but does not specify enough data to recover a complete mathematical statement. The following preserves its proposed assertions and their unresolved qualifications.

In the usual complex-algebraic interpretation, a Fano variety with log-terminal singularities is a normal projective variety \(X\) over \(\C\) with Kawamata log-terminal singularities, with \(K_X\) \(\Q\)-Cartier and \(-K_X\) ample. This interpretation is suggested by the holomorphic terminology; the source fragment does not explicitly specify the ground field or the conventions for families. Write \(\pi\colon\mathcal X\to T\) for a family under consideration and \(X_t\) for its fibre over \(t\in T\). A fibre-tangent holomorphic vector field \(V\) restricts to a vector field \(V_t\) on each fibre. The supplied clean statement assumes that every pair \((X_t,V_t)\) is modified K-semistable. The modification, admissible vector fields, and stability conventions are not defined in the supplied record.

Subject to specifying these missing data, the proposed questions are:
\begin{enumerate}
\item Does the moduli functor of such families admit a proper coarse algebraic moduli space \(M\), parametrizing modified K-polystable Fano objects, such that every connected component of \(M\) is K\"ahler? The functor, its equivalence relation, and the relation between semistable objects and polystable representatives remain unspecified. It is also unclear whether points of \(M\) retain the vector field as part of the object.
\item Do the indicated limiting classes \(L_{T,k}\) exist and define a K\"ahler class, or an ample real line bundle \(L_T\), on \(M\)? Neither \(L_{T,k}\), the meaning of \(T\) in this notation, the limiting process, nor the meaning of an ample real line bundle is supplied. The relationship between the two proposed positivity formulations is therefore unresolved.
\item Given a family of the stated type over a punctured curve, does it admit a modified K-polystable Fano filling, and does such a filling minimize the quantity denoted \(\deg L_T\) among all Fano fillings of that family? Here a filling is intended to extend the given family across the missing point. The allowable extensions, any permitted base changes, extension requirements for \(V\), and the definition of the degree comparison are missing. In particular, a degree attached only to a class on \(M\) does not by itself define the proposed comparison for arbitrary Fano fillings.
\end{enumerate}
The existence of a filling and the minimization assertion cannot be separated unambiguously from the fragment. No complete quantified conjecture is asserted here beyond this explicitly qualified reconstruction.
\subsection{Short English statement}
The supplied problem asks for a proper moduli space of modified K-polystable Fano objects, positive limiting classes on that space, and a degree-minimizing polystable filling of a punctured-curve family. Essential definitions are missing, so this is a qualified statement-only reconstruction requiring mathematical review. The supplied open status and verification labels are curation claims, not independently verified conclusions.
\subsection{Sources}
\begin{itemize}
\item Local UnsolvedMath registry, permanent ID \(30002654\), source problem code \texttt{catalogue entry}, title \emph{K-Moduli and Minimal Fano Fillings}; supplied complete source record from the Oberwolfach Reports Open Problems collection. The project identifies \texttt{data/unsolvedmath/problems.json} and \texttt{catalog/problems.json} as registry resources.
\item Supplied report citation: \emph{Komplexe Analysis (2015)}, Oberwolfach Reports, pp.\ 2177--2234, DOI \texttt{10.4171/owr/2014/39}. Supplied source URL and literature-source link, labelled ``original Oberwolfach report'': \url{https://doi.org/10.4171/owr/2014/39}. The record gives proposed year 2014 and names no proposer. Its background describes the reports as documenting mathematical workshops and their open-problem discussions.
\item Supplied literature assessment, checked exactly \texttt{2026-08-21}: status ``open''; the record is described as an incomplete fragment referring to Odaka-type K-moduli and filling-in conjectures, whose precise asserted statement is not recoverable and whose status was normalized to open. This is supplied dated curation, not independently verified literature assessment. The record itself requests independent verification of time-sensitive status and cited claims.
\item The supplied statement status is \texttt{corrected\_verified}, accompanied by a claim that comparison with the official report restored the hypotheses and notation. This supplied verification claim has not been independently checked and is in tension with the supplied assessment that the precise statement is not recoverable.
\end{itemize}
\end{document}
